\documentclass[usenames,dvipsnames]{beamer}

\mode<presentation> {
\usetheme{Montpellier}
}

\usepackage{graphicx} % Allows including images
\usepackage{booktabs} % Allows the use of \toprule, \midrule and \bottomrule in tables
\usepackage[frenchb]{babel}
\usepackage[T1]{fontenc}
\usepackage{algorithm2e}
\usepackage[utf8]{inputenc}
\setcounter{tocdepth}{2}
\graphicspath{{../Figures/}}


\DeclareMathOperator*{\argmin}{arg\,min}
\DeclareMathOperator*{\argmax}{arg\,max}
\newtheorem*{proposition}{Proposition}

% exercices comptés
\newcounter{numexos}%Création d'un compteur qui s'appelle numexos
\setcounter{numexos}{0}%initialisation du compteur
\newcommand{\exercice}[1]{%Création d'une macro ayant un paramètre
\addtocounter{numexos}{1}%chaque fois que cette macro est appelée, elle ajoute 1 au compteur numexos
\textcolor{DarkOrchid}{Exercice\,\thenumexos\,:}\,#1%Met en rouge Exercice et la valeur du compteur appelée par \thenumeexos
}

\setbeamertemplate{title page}{
    \centering

    \huge\bfseries
    Fondamentaux théoriques du machine learning

\begin{figure}[!h]
\includegraphics[width=8cm]{logistic}
\end{figure}

    }

\title[FTML]{title}


\date{} % Date, can be changed to a custom date

\makeatletter
\let\@@magyar@captionfix\relax
\makeatother

\begin{document}

\begin{frame}
\titlepage % Print the title page as the first slide
\end{frame}

\section{PCA}

\begin{frame}{First principal component}
    We look for $w$, $||w||=1$ such that
    \begin{equation}
	\sum^{n}_{i=1} \big(w^Tx_i \big)^2
    \end{equation}
     is maximal.
    \begin{proposition}
	$w$ is the eigenvector of $X^TX$ with largest eigenvalue $\lambda_{max}$.
    \end{proposition}
\end{frame}

\begin{frame}{First principal component}
    We look for $w$, $||w||=1$ such that
    \begin{equation}
	\sum^{n}_{i=1} \big(w^Tx_i \big)^2
    \end{equation}
     is maximal.
    \begin{proposition}
	\label{prop:pca}
	$w$ is the eigenvector of $X^TX$ with largest eigenvalue $\lambda_{max}$.
    \end{proposition}

    \exercice{\textbf{}} Show the proposition.
\end{frame}

\begin{frame}{First principal component}
        \begin{equation*}
        \begin{aligned}
	    \sum^{n}_{i=1} (w^Tx_i)^2 &= ||Xw||^2\\
	    &= \langle Xw, Xw \rangle\\
	    &= \langle (X^TX)w, w \rangle
        \end{aligned}
    \end{equation*}

    This quantity is always smaller that $\lambda_{max}$, and it attained for an eigenvector in the eigenspace with norm $1$, since we impose that $||w||=1$.

\end{frame}

\end{document} 
